\section{基準・指針・ガイダンス}

\subsection{退職給付会計の範囲}

\begin{InyouJobunbox}
  \begin{itemize}
    \item 基準3
    \item 基準結論50
    \item 指針2
    \item ガイダンス1.1
  \end{itemize}
\end{InyouJobunbox}

% 合スト1-2、3-1

\begin{shortanswer}[年金2 \ 2018 \ 問題2(2)\textcircled{3}]
  役員（取締役、会計参与、監査役および執行役等）の退職慰労金の取扱いについて簡記しなさい。
\end{shortanswer}

\begin{answerbox}
役員の退職慰労金は会計基準の適用範囲に含めない。
なお、厚生年金基金制度及び確定給付企業年金制度の役員部分は会計基準の適用対象となる。その計算にあたって、従業員部分と合わせることができる。
\end{answerbox}

\begin{explanationbox}
前段は基準3に、後段は指針2に記載されている内容。
\end{explanationbox}

\newpage

\subsection{確定拠出制度と確定給付制度}

\begin{InyouJobunbox}
  \begin{itemize}
    \item 基準4〜5、31〜33
    \item 基準結論51、78、78-2、78-3
    \item ガイダンス1.2
  \end{itemize}
\end{InyouJobunbox}

% 合スト1-3、3-2

\begin{shortanswer}[年金2 \ 2022 \ 問題2(6)ア]
  退職給付に関する会計基準第4項に規定されている「確定拠出制度」の定義について簡
潔に入力しなさい。
\end{shortanswer}

\begin{answerbox}
  一定の掛金を外部に積み立て、事業主である企業が当該掛金以外に退職給付に係る追加的な拠出義務を負わない退職給付制度。
\end{answerbox}


\begin{shortanswer}[年金2 \ 2022 \ 問題2(6)イ]
  退職給付に関する会計基準第31項に規定されている「確定拠出制度の会計処理」の内容
について簡潔に入力しなさい。
\end{shortanswer}

\begin{answerbox}
当該制度に基づく要拠出額をもって費用処理する。また、当該制度に基づく要拠出額をもって費用処理するため、未拠出の額は未払金として計上する。
\end{answerbox}

\begin{truefalse}[年金2 \ 2018 \ 問題1(1)\textcircled{2}]
\textbf{下線部の内容の正誤を判定し、誤っている場合は正しい内容に改めなさい。}
\\

「退職給付に関する会計基準」では、確定拠出制度については、当該制度に基づく要拠出額
をもって費用処理する。また、当該制度に基づく要拠出額をもって費用処理するため、未拠
出の額は\underline{翌期の退職給付費用}として計上するものとされている。
\end{truefalse}

\begin{answerbox}
  \textbf{誤}

  未払金
\end{answerbox}

\begin{truefalse}[令和5年第\ajRoman{2}回 \ 短答式公認会計士試験 \ 財務会計論 \ 問題14 \ エ]
\textbf{以下の文章の正誤を判定し、誤っている場合はその理由を記述しなさい。}
\\

確定拠出制度については，当該制度に基づく要拠出額をもって費用処理する。また，
 当該制度に基づく要拠出額をもって費用処理するため，未拠出の額は未払金として計上
 する。
\end{truefalse}


\begin{answerbox}
  \textbf{正}
\end{answerbox}